\begin{figure}[ht]
\centering
\begin{tikzpicture}[>=Stealth, font=\small, node distance=1.05cm]
  \node[draw, rounded corners, fill=blue!8, minimum width=4.8cm, minimum height=0.9cm] (e1) {习题 1：把指数归约到 $4$ 或奇素数};
  \node[draw, rounded corners, fill=green!10, below=of e1, minimum width=4.8cm, minimum height=0.9cm] (e2) {习题 2：用无穷递降解决 $n=4$};
  \node[draw, rounded corners, fill=yellow!16, below=of e2, minimum width=4.8cm, minimum height=0.9cm] (e3) {习题 3：计算 Frey 曲线判别式与最小判别式};
  \node[draw, rounded corners, fill=orange!14, below=of e3, minimum width=4.8cm, minimum height=0.9cm] (e4) {习题 4：验证 $\Sk_2(\GammaO(2))=0$};
  \node[draw, rounded corners, fill=purple!10, below=of e4, minimum width=4.8cm, minimum height=0.9cm] (e5) {习题 5：说明 $X_0(11)\to E_{11}$ 的存在};
  \node[draw, rounded corners, fill=red!8, below=of e5, minimum width=4.8cm, minimum height=0.9cm] (e6) {习题 6：解释为什么是 $3$--$5$ 技巧};

  \draw[->, very thick] (e1) -- (e2);
  \draw[->, very thick] (e2) -- (e3);
  \draw[->, very thick] (e3) -- (e4);
  \draw[->, very thick] (e4) -- (e5);
  \draw[->, very thick] (e5) -- (e6);
\end{tikzpicture}
\caption{终章习题的路线图：从指数归约与无穷递降出发，经 Frey 曲线与空空间 $\Sk_2(\GammaO(2))$，最后回到 Wiles 证明中最精巧的 $3$--$5$ 技巧。}
\label{fig:ch10-exercises-roadmap}
\end{figure}
